\section{本章小结}
\label{sec:transp-summary}

本章从一个问题出发：给定速度场，一滴墨水此后会怎样？答案取决于要问的量。若问粒子携带的值，它沿轨迹不变；若问单位体积的质量，它要除以流动带来的体积倍率。

两种量都通过同一组对象计算。轨迹由常微分方程$\dot{\bm{x}}=\bm{v}(\bm{x},t)$给出，把所有轨迹放在一起得到流映射$\bm{\varphi}_{s,t}$。场值的解是初值与逆流映射的复合，加上沿轨迹累积的源项；密度的解还要多除一个Jacobian行列式$J_{0,t}$，它由速度的散度沿轨迹积分得到。表~\ref{tab:transp-value-density}给出了两者的对比。

这些结论都有明确的条件。

\begin{itemize}
  \item 条件（V）保证了过每一点恰有一条特征线，逆映射在全空间存在。轨迹爆炸时，解只在流映射的值域上有定义；轨迹不唯一时，反向追踪的结果不确定，$C^{1}$解可能不存在。
  \item 总质量守恒需要闭区域、无源项，或全空间上足够的衰减。沿轨迹的密度不变是另一个结论，它只在无散度时成立。
  \item 经典解要求初值$C^{1}$，连续性方程的经典解还要求$\nabla\cdot\bm{v}$对$\bm{x}$连续可微。折角和间断用弱形式处理，弱形式下的存在和唯一仍然要靠条件（V），粗糙速度场下可能出现集中或非唯一。
  \item 数据误差的$L^{p}$放大由体积膨胀控制，速度误差的$L^{\infty}$估计需要初值和源项的Lipschitz常数。间断初值下一致估计失效，只能用$L^{p}$或弱收敛的语言。
\end{itemize}

这一切都有一个共同的前提：速度场是事先给定的，特征线彼此不相交。若传播速度依赖未知量本身，不同高度的点走得快慢不同，特征线就可能相交，解在有限时间内失去光滑性。第~\ref{chap:scalar-conservation-laws}章讨论这种情形。
